\subsection{Right enzyme class, wrong reaction: specificity, promiscuity and missing terms}\label{sec:dd-enzyme}
\textbf{Claim.} MF terms name the right enzyme class but the wrong substrate, chain length, cofactor
or metal, record promiscuous in-vitro activity as function, or have no term for the real reaction.
\textbf{Projects.} \proj{ENZYME\_SPECIFICITY.md}: 14 exemplars incl.\ all ten human mitochondrial
$\beta$-oxidation enzymes (chain-length specificity), LPL1, PHYKPL, EryCIII.
\proj{METEA\_MLL\_CLUSTER.md}: a lanthanide-chelator (lanthanophore) cluster annotated as iron-siderophore
machinery by homology. \proj{RADICAL\_SAM.md}: across 29 members of the largest enzyme superfamily the
problem is resolution, not error (122 modified vs 51 removed). \proj{CYTOCHROME\_P450.md}: 68\% of the
297 experimentally characterised reactions of human P450s reach no GO term.
\proj{NONPHYSIOLOGICAL\_REACTIONS.md}: $\sim$15 MF terms obsoleted for non-physiological substrates
without criteria. \proj{METABOLIC\_MODEL\_ANALYSIS.md}: EC numbers disagree between genome-scale
metabolic models and UniProt for 42--61\% of aligned genes.
\textbf{Rule.} Promiscuity tiers: genuine secondary activity (non-core) / real but of unknown relevance
(over-annotated) / paralog noise when a dedicated enzyme exists (remove, even IDA).
